%% This file contains a number of tweaks that are typically applied to the main document.
%% They are not enabled by default, but can be enabled by uncommenting the relevant lines.

%%
%% Inline annotations; for predefined colors, refer to "dvipsnames" in the xcolor package:
%% https://tinyurl.com/overleaf-colors
%%
\newcommand{\red}[1]{{\color{red}#1}}
\newcommand{\todo}[1]{{\color{red}#1}}
\newcommand{\TODO}[1]{\textbf{\color{red}[TODO: #1]}}
%%
%% disable for camera ready / submission by uncommenting these lines  
%%
% \renewcommand{\TODO}[1]{}
% \renewcommand{\todo}[1]{#1}

%%
%% work harder in optimizing text layout. Typically shrinks text by 1/6 of page, enable
%% it at the very end of the writing process, when you are just above the page limit
%%
% \usepackage{microtype}

%%
%% fine-tune paragraph spacing
%%
% \renewcommand{\paragrapfh}[1]{\vspace{.5em}\noindent\textbf{#1.}}

%%
%% globally adjusts space between figure and caption
%%
% \setlength{\abovecaptionskip}{.5em}


%%
%% Allows "the use of \paper to refer to the project name"
%% with automatic management of space at the end of the word
%%
% \usepackage{xspace}
% \newcommand{\paper}{ProjectName\xspace}

%%
%% Commonly used math definitions
%%
% \DeclareMathOperator*{\argmin}{arg\,min}
% \DeclareMathOperator*{\argmax}{arg\,max}

%%
%% Tigthen underline
%%
% \usepackage{soul}
% \setuldepth{foobar}




% \usepackage{amsmath,amssymb,amsthm}

\usepackage{arydshln}
\usepackage[table]{xcolor}  

\usepackage{epigraph}
\usepackage[normalem]{ulem}
%%%%********************************************************************
% fancy quotes
\definecolor{quotemark}{gray}{0.7}
\makeatletter
\def\fquote{%
    \@ifnextchar[{\fquote@i}{\fquote@i[]}%]
           }%
\def\fquote@i[#1]{%
    \def\tempa{#1}%
    \@ifnextchar[{\fquote@ii}{\fquote@ii[]}%]
                 }%
\def\fquote@ii[#1]{%
    \def\tempb{#1}%
    \@ifnextchar[{\fquote@iii}{\fquote@iii[]}%]
                      }%
\def\fquote@iii[#1]{%
    \def\tempc{#1}%
    \vspace{1em}%
    \noindent%
    \begin{list}{}{%
         \setlength{\leftmargin}{0.05\textwidth}%
         \setlength{\rightmargin}{0.05\textwidth}%
                  }%
         \item[]%
         \begin{picture}(0,0)%
         \put(-8,-5){\makebox(0,0){\scalebox{2}{\textcolor{quotemark}{``}}}}%
         \end{picture}%
         \begingroup\itshape}%
 %%%%********************************************************************
 \def\endfquote{%
 \endgroup\par%
 \makebox[0pt][l]{%
 \hspace{0.27\textwidth}%
 \begin{picture}(0,0)(0,0)%
 \put(60,20){\makebox(0,0){%
 \scalebox{2}{\color{quotemark}''}}}%
 \end{picture}}%
 \ifx\tempa\empty%
 \else%
    \ifx\tempc\empty%
       \hfill\rule{110pt}{0.5pt}\\\mbox{}\hfill\tempa,\ \emph{\tempb}%
   \else%
       \hfill\rule{100pt}{0.5pt}\\\mbox{}\hfill\tempa,\ \emph{\tempb},\ \tempc%
   \fi\fi\par%
   \vspace{0.5em}%
 \end{list}%
 }%
 \makeatother
 %%%%********************************************************************
 

\newcommand{\algopt}{\textsc{\texttt{PIXAR}}\xspace}
\newcommand*\circled[1]{\tikz[baseline=(char.base)]{\node[shape=circle,draw,inner sep=.3pt] (char) {#1};}}

\definecolor{orange}{RGB}{178,92,35}
\definecolor{green1}{RGB}{95,145,51}

\definecolor{red1}{RGB}{197,64,57}
\definecolor{blue1}{RGB}{59,130,220}
\definecolor{green2}{RGB}{82,181,150}
\definecolor{purple1}{RGB}{105,93,223}
\definecolor{orange1}{RGB}{164, 47, 54}
\definecolor{green3}{RGB}{94,145,51}

\definecolor{c_step1}{RGB}{60, 105, 199}
\definecolor{c_step2}{RGB}{159,37,28}
\definecolor{c_step3}{RGB}{107,40,157}

\providecommand{\mpsi}{\boldsymbol{\psi}}
%%%%%%%%%%%%%%%
\providecommand{\lin}[1]{\ensuremath{\left\langle #1 \right\rangle}}
\providecommand{\abs}[1]{\left\lvert#1\right\rvert}
\providecommand{\norm}[1]{\left\lVert#1\right\rVert}

\providecommand{\refLE}[1]{\ensuremath{\stackrel{(\ref{#1})}{\leq}}}
\providecommand{\refEQ}[1]{\ensuremath{\stackrel{(\ref{#1})}{=}}}
\providecommand{\refGE}[1]{\ensuremath{\stackrel{(\ref{#1})}{\geq}}}
\providecommand{\refID}[1]{\ensuremath{\stackrel{(\ref{#1})}{\equiv}}}

\providecommand{\R}{\mathbb{R}} %
\providecommand{\N}{\mathbb{N}} %

\providecommand{\E}{{\mathbb E}}
\providecommand{\E}[1]{{\mathbb E}\left.#1\right. }        %
\providecommand{\Eb}[1]{{\mathbb E}\left[#1\right] }       %
\providecommand{\EE}[2]{{\mathbb E}_{#1}\left.#2\right. }  %
\providecommand{\EEb}[2]{{\mathbb E}_{#1}\left[#2\right] } %
\providecommand{\prob}[1]{{\rm Pr}\left[#1\right] }
\providecommand{\Prob}[2]{{\rm Pr}_{#1}\left[#2\right] }
\providecommand{\P}[1]{{\rm Pr}\left.#1\right. }
\providecommand{\Pb}[1]{{\rm Pr}\left[#1\right] }
\providecommand{\PP}[2]{{\rm Pr}_{#1}\left[#2\right] }
\providecommand{\PPb}[2]{{\rm Pr}_{#1}\left[#2\right] }
\providecommand{\derive}[2]{\frac{{\partial}{#1}}{\partial #2} }  %
\providecommand{\dm}{\mathrm{d}}

\DeclareMathOperator*{\argmin}{arg\,min\,}
\DeclareMathOperator*{\argmax}{arg\,max\,}
\DeclareMathOperator*{\supp}{supp}
\DeclareMathOperator*{\diag}{diag}
\DeclareMathOperator*{\Tr}{Tr}

\providecommand{\0}{\mathbf{0}}
\providecommand{\1}{\mathbf{1}}
\renewcommand{\aa}{\mathbf{a}}
\providecommand{\bb}{\mathbf{b}}
\providecommand{\cc}{\mathbf{c}}
\providecommand{\dd}{\mathbf{d}}
\providecommand{\ee}{\mathbf{e}}
\providecommand{\ff}{\mathbf{f}}
\let\ggg\gg
\renewcommand{\gg}{\mathbf{g}}
\providecommand{\gv}{\mathbf{g}}
\providecommand{\hh}{\mathbf{h}}
\providecommand{\ii}{\mathbf{i}}
\providecommand{\jj}{\mathbf{j}}
\providecommand{\kk}{\mathbf{k}}
\let\lll\ll
\renewcommand{\ll}{\mathbf{l}}
\providecommand{\mm}{\mathbf{m}}
\providecommand{\nn}{\mathbf{n}}
\providecommand{\oo}{\mathbf{o}}
\providecommand{\pp}{\mathbf{p}}
\providecommand{\qq}{\mathbf{q}}
\providecommand{\rr}{\mathbf{r}}
\renewcommand{\ss}{\mathbf{s}}
\providecommand{\tt}{\mathbf{t}}
\providecommand{\uu}{\mathbf{u}}
\providecommand{\vv}{\mathbf{v}}
\providecommand{\ww}{\mathbf{w}}
\providecommand{\xx}{\mathbf{x}}
\providecommand{\yy}{\mathbf{y}}
\providecommand{\zz}{\mathbf{z}}

\providecommand{\mA}{\mathbf{A}}
\providecommand{\mB}{\mathbf{B}}
\providecommand{\mC}{\mathbf{C}}
\providecommand{\mD}{\mathbf{D}}
\providecommand{\mE}{\mathbf{E}}
\providecommand{\mF}{\mathbf{F}}
\providecommand{\mG}{\mathbf{G}}
\providecommand{\mH}{\mathbf{H}}
\providecommand{\mI}{\mathbf{I}}
\providecommand{\mJ}{\mathbf{J}}
\providecommand{\mK}{\mathbf{K}}
\providecommand{\mL}{\mathbf{L}}
\providecommand{\mM}{\mathbf{M}}
\providecommand{\mN}{\mathbf{N}}
\providecommand{\mO}{\mathbf{O}}
\providecommand{\mP}{\mathbf{P}}
\providecommand{\mQ}{\mathbf{Q}}
\providecommand{\mR}{\mathbf{R}}
\providecommand{\mS}{\mathbf{S}}
\providecommand{\mT}{\mathbf{T}}
\providecommand{\mU}{\mathbf{U}}
\providecommand{\mV}{\mathbf{V}}
\providecommand{\mW}{\mathbf{W}}
\providecommand{\mX}{\mathbf{X}}
\providecommand{\mY}{\mathbf{Y}}
\providecommand{\mZ}{\mathbf{Z}}
\providecommand{\mLambda}{\boldsymbol{\Lambda}}
\providecommand{\mOmega}{\boldsymbol{\Omega}}
\providecommand{\mTheta}{\boldsymbol{\Theta}}
\providecommand{\mpi}{\boldsymbol{\pi}}
\providecommand{\malpha}{\boldsymbol{\alpha}}
\providecommand{\meta}{\boldsymbol{\eta}}
\providecommand{\mxi}{\boldsymbol{\xi}}
\providecommand{\mepsilon}{\boldsymbol{\epsilon}}
\providecommand{\mphi}{\boldsymbol{\phi}}
\providecommand{\momega}{\boldsymbol{\omega}}
\providecommand{\mdelta}{\boldsymbol{\delta}}
\providecommand{\mvarsigma}{\boldsymbol{\varsigma}}
\providecommand{\mnu}{\boldsymbol{\nu}}
\providecommand{\mtheta}{\boldsymbol{\theta}}
\providecommand{\mmu}{\boldsymbol{\mu}}
\providecommand{\mf}{\boldsymbol{f}}
\providecommand{\mv}{\boldsymbol{v}}
\providecommand{\mg}{\boldsymbol{g}}
\providecommand{\mmF}{\boldsymbol{F}}
\providecommand{\mmA}{\boldsymbol{A}}
\providecommand{\mmB}{\boldsymbol{B}}
\providecommand{\mmG}{\boldsymbol{G}}

\providecommand{\cA}{\mathcal{A}}
\providecommand{\cB}{\mathcal{B}}
\providecommand{\cC}{\mathcal{C}}
\providecommand{\cD}{\mathcal{D}}
\providecommand{\cE}{\mathcal{E}}
\providecommand{\cF}{\mathcal{F}}
\providecommand{\cG}{\mathcal{G}}
\providecommand{\cH}{\mathcal{H}}
\providecommand{\cI}{\mathcal{I}}
\providecommand{\cJ}{\mathcal{J}}
\providecommand{\cK}{\mathcal{K}}
\providecommand{\cL}{\mathcal{L}}
\providecommand{\cM}{\mathcal{M}}
\providecommand{\cN}{\mathcal{N}}
\providecommand{\cO}{\mathcal{O}}
\providecommand{\cP}{\mathcal{P}}
\providecommand{\cQ}{\mathcal{Q}}
\providecommand{\cR}{\mathcal{R}}
\providecommand{\cS}{\mathcal{S}}
\providecommand{\cT}{\mathcal{T}}
\providecommand{\cU}{\mathcal{U}}
\providecommand{\cV}{\mathcal{V}}
\providecommand{\cX}{\mathcal{X}}
\providecommand{\cY}{\mathcal{Y}}
\providecommand{\cW}{\mathcal{W}}
\providecommand{\cZ}{\mathcal{Z}}

\newenvironment{talign}
{\let\displaystyle\textstyle\align}
{\endalign}
\newenvironment{talign*}
{\let\displaystyle\textstyle\csname align*\endcsname}
{\endalign}



\usepackage[utf8]{inputenc}         %
\usepackage[T1]{fontenc}            %
\usepackage{url}                    %
\usepackage{booktabs}               %
\usepackage{amsfonts}               %
\usepackage{nicefrac}               %
\usepackage{microtype}              %
% \usepackage{xcolor}                 %
\usepackage{algorithm}
\usepackage{algorithmic}
\usepackage{graphicx}
\usepackage{subcaption}
\usepackage[flushleft]{threeparttable}
\usepackage{float}
\usepackage{multirow}
\usepackage{xspace}
% \usepackage{natbib}
\usepackage{enumitem}
\usepackage[font=small]{caption}
\usepackage{autobreak}
\usepackage{sidecap}
\usepackage{wrapfig}
\usepackage{bbding}
\usepackage[toc, page, header]{appendix}
\usepackage{tikz}
% \usepackage{xcolor}
\usepackage{pifont}
\usepackage{mdframed}
\usepackage{colortbl}

% \hypersetup{
%   colorlinks=true,
%   linkcolor=blue,
%   citecolor=blue,
%   urlcolor=blue
% }

\definecolor{coral}{RGB}{255,127,80}
\definecolor{darkgreen}{RGB}{0,100,0}
\definecolor{darkyellow}{RGB}{204,153,0}
\definecolor{salmon}{RGB}{250,128,114}
\definecolor{darkred}{RGB}{150,0,0}
\newcommand{\darkredtext}[1]{{\color{darkred}#1}}

\newcommand{\secref}[1]{\hyperref[#1]{\darkredtext{Sec.~\ref*{#1}}}}
\newcommand{\thmref}[1]{\hyperref[#1]{\darkredtext{Thm.~\ref*{#1}}}}
\newcommand{\defref}[1]{\hyperref[#1]{\darkredtext{Def.~\ref*{#1}}}}
\newcommand{\propref}[1]{\hyperref[#1]{\darkredtext{Prop.~\ref*{#1}}}}
\newcommand{\assumpref}[1]{\hyperref[#1]{\darkredtext{Assump.~\ref*{#1}}}}
\newcommand{\remarkref}[1]{\hyperref[#1]{\darkredtext{Rem.~\ref*{#1}}}}
\newcommand{\hypref}[1]{\hyperref[#1]{\darkredtext{Hyp.~\ref*{#1}}}}
\newcommand{\conjref}[1]{\hyperref[#1]{\darkredtext{Conj.~\ref*{#1}}}}
\newcommand{\lemref}[1]{\hyperref[#1]{\darkredtext{Lem.~\ref*{#1}}}}
\newcommand{\corref}[1]{\hyperref[#1]{\darkredtext{Cor.~\ref*{#1}}}}
\newcommand{\noteref}[1]{\hyperref[#1]{\darkredtext{Nota.~\ref*{#1}}}}
\newcommand{\claimref}[1]{\hyperref[#1]{\darkredtext{Clm.~\ref*{#1}}}}
\newcommand{\obsref}[1]{\hyperref[#1]{\darkredtext{Obs.~\ref*{#1}}}}
\newcommand{\algref}[1]{\hyperref[#1]{\darkredtext{Alg.~\ref*{#1}}}}
\newcommand{\figref}[1]{\hyperref[#1]{\darkredtext{Fig.~\ref*{#1}}}}
\newcommand{\tabref}[1]{\hyperref[#1]{\darkredtext{Tab.~\ref*{#1}}}}
\newcommand{\appref}[1]{\hyperref[#1]{\darkredtext{App.~\ref*{#1}}}}
\renewcommand{\eqref}[1]{\hyperref[#1]{\darkredtext{Eq.~\ref*{#1}}}}



% ============ packages needed for Figure 5 ============
\usepackage{listings}
\usepackage{xcolor}
\usepackage{booktabs}
\usepackage{amssymb}      % for \cup
\usepackage{tikz}         % for \circled (remove if you use the fallback below)

% ============ circled letter marker ============
% preferred version (needs tikz)
\renewcommand{\circled}[1]{%
  \tikz[baseline=(char.base)]{%
    \node[shape=circle,draw,inner sep=1pt,minimum size=1.1em] (char) {\scriptsize #1};}}
% fallback without tikz -- comment out the version above and uncomment this:
% \newcommand{\circled}[1]{\raisebox{.5pt}{\textcircled{\raisebox{-.9pt}{\scriptsize #1}}}}

% ============ listing style for the agent-loop pseudocode ============
\lstdefinestyle{agentloop}{
  basicstyle=\ttfamily\scriptsize,
  commentstyle=\color{gray!70!black}\itshape,
  keywordstyle=\color{blue!55!black}\bfseries,
  showstringspaces=false,
  columns=fullflexible,
  keepspaces=true,
  escapeinside={(*@}{@*)},
  morekeywords={while,if,not,is,continue,and,or,True,False,None,def,return},
  frame=tb,
  framerule=0.5pt,
  rulecolor=\color{gray!60},
  framesep=5pt,
  backgroundcolor=\color{gray!4},
  xleftmargin=8pt,
  xrightmargin=4pt,
  aboveskip=2pt,
  belowskip=2pt,
  breaklines=false,
}


\definecolor{darkred}{RGB}{165,42,42}
\newcommand{\hi}[1]{\textcolor{darkred}{#1}}
